\chapter{What it means to delete a block of labels}
\label{tx:chap:deletion}\label{tx:sec:transfer}
We now describe the operation that the preceding contract requires.
It acts on spaces with label permutations, not by setting some polynomial
variables to zero. Its construction will justify the word ``deletion''
on bases, and then explain why extra work is needed on subspaces.

The construction has three ingredients: spaces on labelled tuples,
one nilpotent operator built from block translations, and a quotient
of its kernel. We develop these tools before stating the exactness
theorem. Its clauses identify the surviving basis, preserve the
actual coefficient maps, and transport the subspace calculation
of a proof. The common-operator construction in the later chapters
establishes these properties.

\section{Why substitution is the wrong operation}
Consider the totality axiom $\sum_jx_{ij}-1$. If the row label $i$ is
deleted and all its variables are set to zero, this polynomial becomes
$-1$. Such a substitution refutes the remaining system immediately.
It does not produce the intended smaller uniform instance.

Restriction of coefficient matrices to retained rows and columns is
correct on uniform data, but arbitrary subspaces do not restrict exactly.
A linear relation supported partly on deleted coordinates may change its
meaning. The proof therefore uses a quotient operation on modules with
label action, together with a theorem that this operation preserves the
required subspace calculations.

\section{Permutation modules and their subspaces}
Recall the label-action language of Section~\ref{tx:sec:action-primer}.
The monomial space is a permutation module, and the NS/PC consequence
spaces are invariant subspaces: relabelling a legal proof gives another
legal proof from the uniform axioms.

\subsection*{A vector space with named basis vectors}
For a finite set of labelled objects $\mathcal U_N$, the notation
$K[\mathcal U_N]$ means the vector space with one basis vector $e_u$
for each object $u\in\mathcal U_N$. A vector is a coefficient list
\[
                         v=\sum_{u\in\mathcal U_N}a_u e_u.
\]
The basis objects might be variables, monomials, axiom indices or
ordered tuples. A label permutation $g\in S_N$ acts by
$g e_u=e_{g u}$ and extends linearly. In matrix language it permutes
the basis coordinates. The phrase \emph{permutation module} means
exactly this vector space together with these permutation matrices.

For instance, the space on single labels has basis
$e_1,\ldots,e_N$, and the swap $(12)$ sends $e_1$ to $e_2$,
$e_2$ to $e_1$, and fixes the other basis vectors. More general
spaces below need not have a basis permuted by labels: an invariant
subspace or a quotient still carries linear permutation actions.
Such a space is called a representation, or module. No additional
algebraic operation on its vectors is intended by that word.

\subsection*{Subspaces, quotients and maps which respect relabelling}
A subspace $U\subseteq M$ is \emph{invariant} if every label
permutation maps it to itself. Its quotient $M/U$ consists of classes
$[v]$ with $v$ and $v'$ representing the same class exactly when
$v-v'\in U$. Relabelling acts by $g[v]=[gv]$, which is well defined
because $gU=U$.

As a small example, in the space on three labels the line
$U=K(e_1+e_2+e_3)$ is invariant. In the quotient the sum
$e_1+e_2+e_3$ becomes zero, while $[e_1]$ is still nonzero.
Thus a quotient does not discard a designated subset of coordinates:
it imposes the specified linear relations. A \emph{subquotient} is
a quotient of an invariant subspace, such as $U/V$ with
$V\subseteq U\subseteq M$ invariant. These are precisely the
objects which arise when we look at targets modulo consequences.

A linear map $f:M\to M'$ is \emph{equivariant} if
$f(gv)=g f(v)$ for every label permutation $g$. This just says that
the two ways of relabelling and applying $f$ give the same answer.
For example, on the space $E_N$ with basis $e_1,\ldots,e_N$,
the coordinate-sum map
\[
 \epsilon:E_N\longrightarrow K,
                 \qquad\epsilon\Bigl(\sum_i a_i e_i\Bigr)=\sum_i a_i
\]
is equivariant when labels act as the identity on the output $K$.
Forgetting one position of a tuple is another example. The axiom,
target and multiplication maps in the uniform proof calculation are
equivariant too; their actual coefficients specify which maps they are.

\subsection*{One expression in permutations, acting on many spaces}
The group algebra $K[S_N]$ consists of finite sums $\sum_g a_g g$.
Its multiplication is composition, extended bilinearly. It acts by the
corresponding linear combination of renamings. If a map
$f$ respects label permutations, then
\[
                           f\bigl((\sum_g a_gg)v\bigr)
                                      = (\sum_g a_gg)f(v).
\]
This elementary identity is the reason a common group-algebra expression
is so valuable: the same expression automatically commutes with every
map in the uniform diagram.

For example, $h=2g-\tau$ acts on $M$ as
$v\mapsto2g v-\tau v$. The same expression acts on $M'$ and
commutes with $f$ whenever $f$ is equivariant. We do not choose
unrelated matrices separately on the two spaces. Nor do we claim
that $h$ commutes with every label permutation: the statement is
that every equivariant map commutes with the common expression.
In positive characteristic, the coefficient $2$ in this example is
interpreted in $K$, so it is zero when the characteristic is two.

\subsection*{Why ordered tuples are the test spaces}
For $b\ge0$ let
\[
                 P_b(N)=K[\Inj([b],[N])]
\]
be the space on ordered tuples of distinct labels. Every orbit of a
support-$b$ index is a quotient of such a space. Choose an ordering of
its support; an injective tuple specifies where that support is sent.
Two extensions agreeing on the support act identically on the index.
Thus ordered tuples provide a finite universal list of spaces on which
to construct an operator valid for all support-bounded data.
Here $P_b(N)$ denotes a tuple permutation space; the bounded-degree
polynomial space is written $P_{N,\le d}$.

Here $[b]=\{1,\ldots,b\}$ names the \emph{positions}; an injection
assigns distinct universe labels to them. In particular:
\begin{itemize}
\item $P_0(N)$ is the one-dimensional space on the empty tuple.
Every permutation fixes its basis vector.
\item $P_1(N)$ has basis $e_1,\ldots,e_N$.
\item $P_2(N)$ has basis $e_{(i,j)}$ for $i\ne j$. It has
$N(N-1)$ basis vectors, and $g e_{(i,j)}=e_{(g(i),g(j))}$.
\end{itemize}
There is also an action permuting tuple \emph{positions}: swapping
the two positions sends $(i,j)$ to $(j,i)$. It is different from
swapping two particular universe labels. The two actions commute.
Throughout the construction, $b$ counts positions and $r$ bounds
their number; neither is a polynomial degree.

The universality claim in the preceding paragraph is concrete. If an
object $u$ has support $\{u_1,\ldots,u_b\}$, send
$e_{(t_1,\ldots,t_b)}$ to $e_{g u}$, where $g$ sends each $u_i$
to $t_i$. Any two such permutations agree on the support and hence
give the same object. The map is linear, equivariant and onto the
space on the orbit of $u$. Therefore any operator which vanishes on
$P_b(N)$ also vanishes on that orbit space.

\subsection*{The common-annihilator class}
Let $J_r(N)$ be the ideal of group-algebra elements acting as zero on
$P_b(N)$ for every $b\le r$. Write $\mathcal C_r(N)$ for the
finite-dimensional modules annihilated by $J_r(N)$. Every support-$r$
permutation basis, and every submodule or quotient of it, belongs to this
class. It is also closed under finite direct sums. We will not assume
that it is closed under every tensor product: the actual tensor products
used in the proof are placed inside known support-bounded permutation
ambients.

In formulas, the definition is
\[
 J_r(N)=\{h\in K[S_N]:h v=0\text{ for every }v\in P_b(N)
                                  \text{ and every }0\le b\le r\}.
\]
Calling it an \emph{ideal} records that sums and scalar multiples
remain in it, as do products on either side by group-algebra
elements. The latter fact follows because each $P_b(N)$ is preserved
by all those operators.

Membership of $M$ in $\mathcal C_r(N)$ means that \emph{every}
operator in this set acts as zero on $M$ too. Equivalently, all the
operator relations shared by the tuple spaces up to length $r$ hold
on $M$. This is a condition on the permutation action on a space,
not a claim that each of its vectors mentions only $r$ labels.
For example, $\sum_{i=1}^N a_i e_i$ can use every label, while its
single-label permutation space belongs to $\mathcal C_1(N)$.
Likewise the degree-$d$ monomial space has basis support at most
$kd$, although a polynomial in it can involve many different labels
across its monomials.
The symbol $J_r(N)$ denotes this set of operators, not the target
map $J_N$ from the previous chapter.

The closure assertions now have elementary proofs. If an operator
kills every vector of $M$, it kills every vector of an invariant
subspace and acts as zero on every quotient class. It also kills
a finite direct sum of such spaces. Hence the class includes the
subspaces and subquotients needed by a proof. It need not include
an arbitrary tensor product: combining two basis objects can add
their support sizes. We will check the relevant \emph{ambient}
support each time a tensor product is used.

\section{The regular prime-power block}
The parameters have different roles. We retain $n$ labels, add
$q=p^a$ block labels, and obtain $N=n+q$ universe labels in total.
Here $p$ is the characteristic, $a$ is a block exponent, and $r$
is the support budget for the spaces in the proof diagram. Choosing
$q>r$ does not mean choosing $q$ larger than $n$.

Write $N=n+q$ and identify the last $q=p^a$ labels with
$A=(\mathbb F_p)^a$. Translation by a coordinate vector gives a
permutation $\sigma_i$ of the block, fixing the retained labels. The
translations commute and have order $p$. Over a suitable coefficient
extension choose $c_1,\ldots,c_a$ independent over $\mathbb F_p$ and put
\[
                 z=\sum_{i=1}^a c_i(\sigma_i-1).
\]
Here $1$ is the identity permutation. On a representation $M$, this
group-algebra element acts as a linear operator $z_M$. A power of $z$
means repeated composition: $z^p=0$ will say that $p$ successive
applications send every vector to zero.
The characteristic-$p$ binomial formula gives $z^p=0$. The retained
permutations commute with $z$, so its kernel and the quotient below
retain an action of $S_n$.
In operator formulas, $I$ denotes the identity linear map, the action
of this identity permutation. Independence of the coefficients is
over the prime field $\mathbb F_p$: a relation
$\sum_i t_i c_i=0$ with $t_i\in\mathbb F_p$ must have all
$t_i=0$. It is not independence over the larger working field.

The word \emph{regular} means that the translations move any block
label to any other in exactly one way. For $p=2,a=2$, write the
four block labels as $00,10,01,11$. Translation in the first
coordinate swaps $00$ with $10$ and $01$ with $11$; translation
in the second swaps $00$ with $01$ and $10$ with $11$. They
leave all retained labels fixed. In general, the identity
\[
 z^p=\sum_i c_i^p(\sigma_i-I)^p
     =\sum_i c_i^p(\sigma_i^p-I)=0
\]
uses the commutation of these translations and the characteristic-$p$
binomial formula. It is a statement about matrices or composed
operators, not about a scalar whose power is zero.

The independence of the coefficients prevents a nontrivial translation
orbit from being accidentally invisible. It can require extending the
field. Membership of a target with original-field coefficients descends
by Lemma~\ref{tx:lem:field}; no proof over an arbitrary extension is
silently substituted for a proof over the original field.

\subsection*{A quotient made from a kernel and an image}
Define
\[
                 \Del(M)=\ker z_M/(\ker z_M\cap\im z_M).
\]
For $p=2$ this is $\ker z/\im z$. For odd $p$ that shorter formula
is not meaningful in general, because $z^2$ need not be zero.

The kernel consists of vectors $v$ with $z_Mv=0$. The image
consists of vectors $z_Mw$. We first keep the kernel vectors and
then identify two of them when their difference is also an image
vector. Thus an element of $\Del(M)$ is a class $[v]$ of a vector
already in the kernel. An arbitrary basis vector need not define
such a class.

For a fully explicit case, work in characteristic two, retain one
label $a$, and swap two block labels $b,c$. Choose
$z=(bc)-I$ on the space with basis $e_a,e_b,e_c$. Then
\[
 ze_a=0,\qquad ze_b=e_b+e_c,\qquad ze_c=e_b+e_c.
\]
Consequently
\[
 \ker z=K e_a\oplus K(e_b+e_c),\qquad
 \im z=K(e_b+e_c),\qquad \Del(M)\simeq K e_a.
\]
The class of $\alpha e_a+\beta(e_b+e_c)$ is represented by
$\alpha e_a$. The combination $e_b+e_c$ lies in the kernel,
so it qualifies as a representative; it also lies in the image,
so it represents the zero class.
The vector $e_a$ survives. The individual vector $e_b$ is not in
the kernel and hence is not a representative for this quotient.
This distinction is why deletion is not a coordinate restriction
applied to every vector.

To see what the quotient keeps, take a nilpotent chain with basis
$v_0,\ldots,v_{\ell-1}$ and
\[
                  zv_0=0,\qquad zv_j=v_{j-1}\quad(j>0).
\]
This is a nilpotent \emph{Jordan block}: applying $z$ moves each vector
one position down the chain and eventually sends it to zero.
The chain basis describes the linear operator $z$. It need not be a
permutation basis or an invariant decomposition for the full symmetric
group. This distinction will matter when we later build an operator
which works on many representations at once.

For $\ell=1$, the kernel is the whole line and the image is zero,
so the line survives. For $\ell\ge2$, the kernel is $Kv_0$ and
$v_0=zv_1$ already lies in the image, so the quotient is zero.
Thus deletion retains length-one chains and discards every longer chain.
The equation $z^p=0$ permits every length $1\le\ell\le p$.
The later odd-characteristic construction separately proves that the
relevant tuple generators have only lengths one and $p$; neither that
purity statement nor exactness follows from nilpotence alone.

\subsection*{How a linear map induces a map of the quotients}
Let $f:M\to M'$ commute with label permutations. The common
group-algebra definition gives $f z_M=z_{M'}f$. If $v\in\ker z_M$,
then $f(v)\in\ker z_{M'}$. If two kernel vectors differ by
$z_Mw$, their images differ by $z_{M'}f(w)$, which is again an
image vector and remains in the kernel. Therefore the rule
\[
                  \Del(f):\Del(M)\longrightarrow\Del(M'),
                              \qquad[v]\longmapsto[f(v)]
\]
is well defined. It is the actual induced map, not an independently
chosen map between spaces of the same dimension. Applying two maps
in succession before or after taking classes gives the same answer.

In the three-label example, the coordinate-sum map
$\epsilon:M\to K$ is equivariant. On a kernel vector it gives
\[
 \epsilon\bigl(\alpha e_a+\beta(e_b+e_c)\bigr)
                       =\alpha+\beta+\beta=\alpha.
\]
Thus its induced map sends the surviving class $[e_a]$ to $1$
and the zero class $[e_b+e_c]$ to $0$. On the trivial output
space $K$ every label permutation is the identity, so $z=0$ and
$\Del(K)=K$.

Finally, permutations of the retained labels commute with each
block translation. They preserve the kernel and image and act on
the quotient by $[v]\mapsto[g v]$. This is why the output carries
an action of the \emph{smaller} group $S_n$. The deleted block
does not become a set of new polynomial variables in that output.

\section{Why small supports make a basis object retained}
The support bound in the next argument belongs to an individual
\emph{basis object}, not to every vector in its span. A basis variable
$x_i$ uses label $i$; the monomial $x_i x_j$ uses labels $i,j$;
a sum of many such monomials may mention most of the universe.

Suppose $\mathcal U$ has support bound $u<q$. If an object is fixed by
all block translations, its support is a block-translation-invariant set.
The regular action on the block is transitive. Hence a support meeting
the block would contain all $q$ block labels, contrary to $u<q$.
The fixed objects are exactly those supported on retained labels.

This is the first structural use of the inequality $q>r$. It has nothing
to do with choosing $q$ larger than the universe itself. It must exceed
the support needed by the degree-$d$ calculation. A one-label change
usually has $q=1\le r$ and cannot satisfy this argument. This explains
why the natural shifts are prime powers rather than $n\mapsto n+1$.

The strict inequality is visible in an example. With a four-label
block and unary variables, the product of all four block variables
is a basis monomial fixed by every translation. It uses the whole
block, so it is excluded by a support bound $u<4$. A monomial
meeting the block in one, two or three labels cannot be fixed by
all translations.
The support argument identifies which basis objects are fixed;
the later orbit calculation must still prove that the moved part
has zero deletion quotient.

\subsection*{The theorem in plain language}
Choose a block larger than the supports we need to track, and keep
enough labels outside it. The theorem asserts that the block
coefficients can be chosen so that the quotient just defined does
three things at once. It retains the expected smaller basis; it
carries the actual uniform coefficient maps to the smaller maps;
and, on the class $\mathcal C_r$, it preserves the subspace
operations used by a proof. The last assertion is the substantial
one: neither the definition of the quotient nor $z^p=0$ proves it.

Here is the precise statement. The phrases ``exact'', ``natural''
and ``trivial module'' are read clause by clause immediately below.

\begin{theorem}[The exact deletion input]\label{tx:thm:deletion}
Work first over an algebraically closed field of characteristic $p$.
For $q=p^a>r$ and retained universe size $n\ge\max\{0,2r-1\}$, the block
coefficients can be chosen so that $\Del$ is exact on $\mathcal C_r(n+q)$.
It fixes the trivial module. For every support-$u$ permutation basis with
$u<q$ there are natural identifications
\begin{align}
 \Del(K[\mathcal U_{n+q}])&\simeq K[\mathcal U_n],\label{tx:eq:permdelete}\\
 \Del(K[\mathcal U_{n+q}]\otimes Y)&\simeq
                  K[\mathcal U_n]\otimes\Del(Y),\label{tx:eq:tensordelete}
\end{align}
with a permutation module as the first tensor factor. Stable invariant
coefficient matrices are carried to their actual retained-universe matrices.
\end{theorem}

This finite-algebra theorem is proved by the chapters on
square-zero blocks, linear layer repair, rook realisation and top-power
projection. No representation-theoretic equivalence is
invoked as an external assumption.

\subsection*{Reading the clauses}
\paragraph{The working field.}
An algebraically closed field has the coefficients needed for the
independent block directions. This is where the construction is
carried out, not a change of the original input equations.
Lemma~\ref{tx:lem:field} brings membership of original-coefficient
targets back to the original field. The block size $q=p^a$ is a
number of universe labels, not a prescribed cardinality for that
field.

\paragraph{The support budget and the range.}
The integer $r$ bounds the tuple test spaces and the supports of the
proof-diagram ambients which must undergo exact transport.
The inequalities are $q>r$ and $n\ge\max\{0,2r-1\}$.
They concern block labels and retained labels, respectively. For
$r\ge1$, the latter is the sufficient retained-label budget
$2r-1$; for $r=0$, it is simply $n\ge0$. The proof will construct
one compatible witness for all lengths $0,\ldots,r$ on that one
retained pool. Template stabilisation and target-degree thresholds
are additional requirements when this finite-algebra input is
applied to polynomial families.

\paragraph{Exactness on the specified class.}
Take an equivariant map $f:M\to M'$ between spaces in
$\mathcal C_r(N)$. Its kernel and image are invariant subspaces.
Exactness says that their deleted spaces are the actual kernel and
image of $\Del(f)$, with the induced inclusions. In particular,
if $U\subseteq V$ are invariant spaces in this class, the induced
map $\Del(U)\to\Del(V)$ is injective, and
\[
                  \Del(V/U)\simeq\Del(V)/\Del(U).
\]
This is a claim about maps and their linear relations, not only
about dimensions. It is restricted to $\mathcal C_r(N)$;
arbitrary label representations are not being declared exact.

\paragraph{Fixing the trivial module.}
On the one-dimensional space $K$ with all label permutations acting
as the identity, every $\sigma_i-I$ acts as zero. Hence $z=0$,
the kernel is $K$, the image is zero and $\Del(K)=K$.
The constant line does not disappear. This clause will be needed
both for invariant targets and for the detection pairing.

\paragraph{The first basis identification.}
The first displayed isomorphism says that the surviving directions
in the space on $\mathcal U_{n+q}$ are precisely the basis
objects supported on retained labels. A retained object $u$ gives
the class of its basis vector $e_u$, identified with the same
object in $\mathcal U_n$. Moved basis vectors need not themselves
be kernel vectors; the assertion is that their entire moved-orbit
summand contributes zero to the quotient. For a single-label basis,
this turns the large label space into the space on the $n$
retained labels.

\paragraph{The tensor identification and its naturality.}
Here $Y$ is any space with the larger label action. The first
factor has a permutation basis. Regard a tensor
$\sum_u e_u\otimes y_u$ as a list of vectors from $Y$, one
packet for each basis object $u$. A label permutation acts on
\emph{both} the packet index and its contents:
$g(e_u\otimes y)=e_{gu}\otimes gy$.
The theorem keeps retained packet indices and applies deletion
to their contents. More precisely, for retained $u$ and
$y\in\ker z_Y$, it matches
\[
                         e_u\otimes[y]\quad\text{with}\quad[e_u\otimes y].
\]
It also proves that the moved packet summand disappears even
though translations act on its $Y$ contents. There is no separate
bounded-support assumption on $Y$ in this basis identification.
When exactness is used on a tensor source, we separately check
that source's ambient support budget.

``Natural'' means that applying any equivariant map $Y\to Y'$
to the packet contents before or after this identification gives
the same result. This is what lets the rule apply to a consequence
subspace $Y\subseteq L_d$ with its actual inclusion, rather than
to a separately chosen basis for $Y$. It is also why the tensor
clause specifically states which factor has a permutation basis.

\paragraph{The coefficient matrices are the actual maps.}
If $F_N$ sends a retained input basis vector $e_u$ to
$\sum_v a_N(v,u)e_v$, equivariance puts this image in the kernel
of $z$. Its retained coefficients are exactly $a_N(v,u)$ for
retained $v$. A stable rule gives $a_N(v,u)=a_n(v,u)$ on those
objects, and the moved kernel component vanishes in the quotient.
Thus the induced map is the prescribed $F_n$ entry by entry.
The statement does not allow us to replace it by an unrelated
matrix with the same rank. It applies in particular to the
uniform axiom, target and multiplication maps from Chapter~\ref{tx:chap:contract}.

\subsection*{How the later proof makes exactness possible}
The later chapters construct one group-algebra element $w$ with
\[
             z^{p-1}wz^{p-1}=z^{p-1}
                           \quad\text{on every }P_b(N),\ b\le r.
\]
This is a partial-inverse condition: on the image of $z^{p-1}$,
$w$ supplies a preimage. It does not invert the nilpotent operator
on the whole space. From $z$ and $w$ the proof constructs another
common operator $e$, using also the proved chain property of the
tuple generators. This $e$ is a projector with $e^2=e$, and identifies
$\Del(M)$ with its image $eM$ as a vector space. The letter $e$
here names an operator, unlike the basis-vector notation $e_u$.

Every actual equivariant map commutes with this same expression
in permutations. For example, if $f:V\to W$ is onto and
$y\in eW$, choose an ordinary preimage $v$ of $y$. Then
$f(ev)=e f(v)=ey=y$, so the restricted map $eV\to eW$
is still onto. The kernel and injection checks use the same
commutation. This is the source of exactness. Its hard part is
constructing one such expression which works on the entire
annihilator class; separate projectors on separate spaces would
not justify these equations. The quotient itself carries the
retained label action; the projector image is its vector-space
model.

\subsection*{An admissible numerical instance}
Take characteristic two, $r=2$, $q=4$, $n=3$, and $N=7$.
Work in an algebraic closure of $\mathbb F_2$, choose
$\omega^2+\omega+1=0$, and use coefficients $1,\omega$.
Retain labels $1,2,3$ and identify labels $4,5,6,7$ with
$00,10,01,11$, respectively. The two translations and their
common operator are
\[
 \sigma_1=(45)(67),\qquad\sigma_2=(46)(57),\qquad
 z=(\sigma_1-I)+\omega(\sigma_2-I).
\]
The numbers satisfy $q>r$ and $n=2r-1$.
The tuple test spaces are $P_0(7)$, $P_1(7)$ and $P_2(7)$,
of dimensions $1,7,42$. The moved-orbit calculation, which uses
the independence of $1,\omega$ over $\mathbb F_2$, identifies
their deletion quotients with retained tuple spaces of dimensions
$1,3,6$. This illustrates the basis calculation for the displayed
operator.

The all-class exactness assertion is separate. For these same
parameters, Theorem~\ref{tx:thm:deletion} provides a suitable choice
of block coefficients for which deletion is exact on
$\mathcal C_2(7)$. The present numerical illustration does not
certify that its specific choice $1,\omega$ has that additional
property or supply a common witness for it. The later construction
establishes the exactness input with its required coefficient choice.

For comparison with polynomial notation, the ordinary degree-at-most-two
monomial space in seven unary variables has $36$ basis monomials:
one constant, seven linear monomials, seven squares and twenty-one
products of distinct variables. Each monomial uses at most two
labels. Its deletion is the corresponding space in three unary
variables, with $10$ monomials. This is a consequence of the basis
identification, not a computation of any minimum NS or PC degree.
Squares are separate coordinates because this is the ordinary
polynomial ambient, not the space of functions on Boolean assignments.
An application to a proof must additionally cover the actual axiom
and target maps, multiplication domains and, when needed, the
doubled-support detector. Those data may require a larger $r$.

The instance illustrates the theorem's parameters and conclusions;
its concrete operator illustrates basis deletion, and the theorem's
existential coefficient choice supplies exactness. It does not
independently construct a common witness or prove general lifting
by finite experiment.

\section{Why deletion on bases is not enough}
In characteristic two, let two block labels be swapped. Their permutation
space has basis $e_1,e_2$ and $z e_1=z e_2=e_1+e_2$. Its kernel equals
its image, so its deletion is zero. The invariant line
$T=K(e_1+e_2)$ has $z=0$ and survives. Therefore deletion of
$T\hookrightarrow Ke_1\oplus Ke_2$ is the zero map on a nonzero space.
An injection was not preserved.

Here there are no retained labels: $n=0$, $q=2$ and the single-label
basis has support budget $r=1$. The retained-universe hypothesis would
require $n\ge1$. The example therefore explains the need for a
range restriction rather than contradicting Theorem~\ref{tx:thm:deletion}.
The vector $e_1+e_2$ represents a nonzero quotient class when
the operation is applied to $T$ alone, but represents the zero
class after inclusion into the larger two-point space. Inducing
maps of quotients is always possible; preserving injectivity is
an additional theorem.

The extra retained labels in Theorem~\ref{tx:thm:deletion} rule out
this problem in the relevant class by providing a common projector.
The basis orbit calculation and the exactness construction have distinct
jobs. Forgetting the second job would leave the proof unable to handle
intersections and hence ordinary-PC degree falls.

\section{What exactness will buy}
In the contract of Definition~\ref{tx:def:polynomial-diagram}, matrix
naturality identifies the actual axiom, multiplication and target maps.
Exactness preserves the kernels, images, intersections and sums of the
closure calculation. The doubled-support pairing supplies the remaining
target-vanishing reflection for general indexed targets.
The roles should remain distinct. Exactness carries a consequence
space and its actual inclusion to the smaller universe. The target-map
compatibility then identifies the deleted request. To reflect
underivability, we still need to know that a nonzero target class
in the relevant quotient cannot disappear. That last step is the
nonzero-subquotient detector, not a consequence of the word
``exact'' alone.

The next chapter gives the related upward lift and explains how it
interacts with multiplication without pretending to be a ring map.
